% ============================================================================
% LANGENTWURF LE-06d: Wiederholung + Stundenleistung
% Spanisch Klasse 7 - Sequenz 6: Mi ciudad
% ============================================================================
% Tier: 1 (vollständig)
% Doppelstunde: 6d (Std. 7-8 von 8) - ABSCHLUSS
% Fachfarbe: #c41e3a (Karminrot)
% ============================================================================

\documentclass[11pt,a4paper]{article}

\usepackage[utf8]{inputenc}
\usepackage[T1]{fontenc}
\usepackage[ngerman]{babel}
\usepackage{geometry}
\usepackage{fancyhdr}
\usepackage{lastpage}
\usepackage{xcolor}
\usepackage{tcolorbox}
\usepackage{enumitem}
\usepackage{tabularx}
\usepackage{longtable}
\usepackage{array}
\usepackage{booktabs}
\usepackage{amssymb}

% ============================================================================
% KONFIGURATION
% ============================================================================

\newcommand{\fach}{Spanisch}
\newcommand{\fachfarbe}{c41e3a}
\newcommand{\jahrgangsstufe}{7}
\newcommand{\schulart}{Gesamtschule}
\newcommand{\stundenthema}{Wiederholung + Stundenleistung}
\newcommand{\reihenthema}{Mi ciudad}
\newcommand{\stundennummer}{6d}
\newcommand{\reihenstunden}{4}
\newcommand{\gession}{A1}
\newcommand{\lehrwerk}{Qué pasa 1}
\newcommand{\lehrwerkseite}{S. 102-119}
\newcommand{\lerngruppengroesse}{24}

% ============================================================================
% LAYOUT
% ============================================================================
\geometry{left=2.5cm, right=2.5cm, top=2.5cm, bottom=2.5cm}
\definecolor{fach}{HTML}{\fachfarbe}
\definecolor{gray}{HTML}{666666}
\definecolor{lightgray}{HTML}{f5f5f5}
\definecolor{successgreen}{HTML}{2a7a4b}
\definecolor{warningorange}{HTML}{b35c00}
\definecolor{basis}{HTML}{2a7a4b}
\definecolor{standard}{HTML}{2c5aa0}
\definecolor{erweiterung}{HTML}{7b1fa2}

\tcbuselibrary{skins,breakable}

\newtcolorbox{infobox}[1][]{
    colback=lightgray,
    colframe=fach,
    fonttitle=\bfseries,
    title=#1,
    breakable
}

\newtcolorbox{scorebox}{
    colback=white,
    colframe=fach,
    fonttitle=\bfseries,
    title=Didaktik-Score,
    breakable
}

\pagestyle{fancy}
\fancyhf{}
\fancyhead[L]{\small\textcolor{gray}{\fach{} -- Jg. \jahrgangsstufe{} -- \stundenthema}}
\fancyhead[R]{\small\textcolor{gray}{LE-\stundennummer{} | Tier 1}}
\fancyfoot[C]{\small\thepage{} / \pageref{LastPage}}
\renewcommand{\headrulewidth}{0.4pt}

% Differenzierungsmarker
\newcommand{\B}{\textcolor{basis}{\textbf{[B]}}}
\newcommand{\Std}{\textcolor{standard}{\textbf{[S]}}}
\newcommand{\E}{\textcolor{erweiterung}{\textbf{[E]}}}

% ============================================================================
% DOKUMENT
% ============================================================================
\begin{document}

% ============================================================================
% DECKBLATT
% ============================================================================
\begin{titlepage}
    \centering
    \vspace*{2cm}
    {\large\textcolor{gray}{\schulart}}\\[2cm]
    {\Huge\textcolor{fach}{\textbf{Langentwurf}}}\\[0.5cm]
    {\large LE-\stundennummer{} | Tier 1 (vollständig)}\\[2cm]
    {\LARGE\textbf{\stundenthema}}\\[0.5cm]
    {\large Sequenz: \reihenthema{} -- \textbf{ABSCHLUSS}}\\[2cm]
    \begin{tabular}{rl}
        \textbf{Fach:} & \fach{} (A1) \\[0.3cm]
        \textbf{Klasse:} & \jahrgangsstufe \\[0.3cm]
        \textbf{Lehrwerk:} & \lehrwerk, \lehrwerkseite \\[0.3cm]
        \textbf{Doppelstunde:} & \stundennummer{} (Std. 7-8 von 8) \\[0.3cm]
    \end{tabular}
    \vfill
\end{titlepage}

\tableofcontents
\newpage

% ============================================================================
% 1. BEDINGUNGSANALYSE
% ============================================================================
\section{Bedingungsanalyse}

\subsection{Lerngruppe}
\begin{tabular}{ll}
    Kursgröße: & \lerngruppengroesse{} SuS \\
    Jahrgangsstufe: & \jahrgangsstufe \\
    Sprachniveau: & A1 (GER) -- Anfänger \\
    Fremdsprache: & 2. Fremdsprache (nach Englisch) \\
\end{tabular}

\subsection{Vorwissen aus 06a/06b/06c}
\begin{itemize}
    \item \textbf{06a:} Orte in der Stadt (la tienda, el parque, el cine, la plaza...)
    \item \textbf{06b:} Verb \textit{ir} (gehen/fahren) + Präposition \textit{a}
    \item \textbf{06c:} Verb \textit{estar} (sich befinden) + Wegbeschreibungen (todo recto, a la derecha/izquierda, cerca de, lejos de)
\end{itemize}

\subsection{Differenzierung (intern)}
\begin{tabular}{|l|p{3.5cm}|p{3.5cm}|p{3.5cm}|}
\hline
& \B{} \textbf{Basis} & \Std{} \textbf{Standard} & \E{} \textbf{Erweiterung} \\
\hline
\textbf{Ziel} & BR: Orte benennen, einfache Sätze & MR: ir/estar anwenden & GY: Wegbeschreibung produzieren \\
\hline
\textbf{Orte} & Zuordnung Bild-Wort & Sätze mit estar bilden & Beschreibung verfassen \\
\hline
\textbf{Test} & Basisaufgaben (60\%) & Alle Pflichtaufgaben & + Zusatzaufgaben \\
\hline
\end{tabular}

% ============================================================================
% 2. SACHANALYSE
% ============================================================================
\section{Sachanalyse}

\subsection{Orte in der Stadt (lugares en la ciudad)}
\begin{center}
\begin{tabular}{|l|l||l|l|}
\hline
\textbf{Spanisch} & \textbf{Deutsch} & \textbf{Spanisch} & \textbf{Deutsch} \\
\hline
la tienda & das Geschäft & el supermercado & der Supermarkt \\
\hline
el parque & der Park & la farmacia & die Apotheke \\
\hline
el cine & das Kino & el hospital & das Krankenhaus \\
\hline
la plaza & der Platz & la estación & der Bahnhof \\
\hline
el museo & das Museum & el banco & die Bank \\
\hline
la iglesia & die Kirche & el restaurante & das Restaurant \\
\hline
\end{tabular}
\end{center}

\subsection{Verben: ir und estar}

\textbf{ir (gehen/fahren):}
\begin{center}
\begin{tabular}{|l|l|l|}
\hline
\textbf{Person} & \textbf{ir} & \textbf{Beispiel} \\
\hline
yo & voy & Voy al cine. \\
\hline
tú & vas & ¿Vas al parque? \\
\hline
él/ella & va & Ella va a la tienda. \\
\hline
nosotros & vamos & Vamos al museo. \\
\hline
\end{tabular}
\end{center}

\textbf{Wichtig:} a + el = \textbf{al}

\vspace{1em}

\textbf{estar (sich befinden):}
\begin{center}
\begin{tabular}{|l|l|l|}
\hline
\textbf{Person} & \textbf{estar} & \textbf{Beispiel} \\
\hline
yo & estoy & Estoy en el parque. \\
\hline
tú & estás & ¿Dónde estás? \\
\hline
él/ella & está & El banco está aquí. \\
\hline
nosotros & estamos & Estamos en la plaza. \\
\hline
\end{tabular}
\end{center}

\subsection{Wegbeschreibungen}
\begin{center}
\begin{tabular}{|l|l|}
\hline
\textbf{Spanisch} & \textbf{Deutsch} \\
\hline
todo recto & geradeaus \\
\hline
a la derecha & nach rechts \\
\hline
a la izquierda & nach links \\
\hline
cerca de & in der Nähe von \\
\hline
lejos de & weit weg von \\
\hline
al lado de & neben \\
\hline
enfrente de & gegenüber von \\
\hline
\end{tabular}
\end{center}

% ============================================================================
% 3. DIDAKTISCHE ANALYSE
% ============================================================================
\section{Didaktische Analyse}

\subsection{Stellung in der Sequenz}
Dies ist die \textbf{4. und letzte Doppelstunde} (Std. 7-8 von 8) der Sequenz ``\reihenthema''.
Die Stunde schließt die Sequenz ab mit:
\begin{enumerate}
    \item Konsolidierung aller Inhalte (Orte, ir, estar, Wegbeschreibung)
    \item Interaktiver Wiederholung mit Lernpfad
    \item Stundenleistung (summative Bewertung)
\end{enumerate}

\subsection{Kompetenzziele (Rahmenplan MV)}
\begin{itemize}
    \item \textbf{Sprechen:} Nach dem Weg fragen und Auskunft geben
    \item \textbf{Verfügen über sprachliche Mittel:} Ortsvokabular, Verben ir/estar
    \item \textbf{Interkulturell:} Orientierung in einer spanischen Stadt
\end{itemize}

\subsection{Legitimation}
\textbf{Rahmenplan Spanisch MV:}
\begin{itemize}
    \item An Gesprächen teilnehmen: einfache Wegbeschreibungen
    \item Verfügen über sprachliche Mittel: Ortsangaben und Richtungsangaben
\end{itemize}

% ============================================================================
% 4. STUNDENZIELE
% ============================================================================
\section{Stundenziele}

\subsection{Grobziel}
Die SuS \textbf{konsolidieren} ihre Kenntnisse zu Orten in der Stadt, den Verben \textit{ir} und \textit{estar} sowie Wegbeschreibungen und weisen ihre Kompetenzen in einer \textbf{Stundenleistung} (20 Punkte) nach.

\subsection{Feinziele (operationalisiert)}
\begin{tabular}{|c|p{8cm}|c|c|}
\hline
\textbf{Nr.} & \textbf{Die SuS können...} & \textbf{AFB} & \textbf{Diff.} \\
\hline
FZ1 & mindestens 10 Orte in der Stadt auf Spanisch benennen und zuordnen & I & alle \\
\hline
FZ2 & das Verb \textit{ir} korrekt konjugieren und mit \textit{a + Artikel} verwenden & II & alle \\
\hline
FZ3 & das Verb \textit{estar} zur Ortsangabe verwenden und konjugieren & II & alle \\
\hline
FZ4 & eine einfache Wegbeschreibung verstehen und produzieren & II/III & \Std{}/\E{} \\
\hline
\end{tabular}

% ============================================================================
% 5. METHODISCHE ANALYSE
% ============================================================================
\section{Methodische Analyse}

\subsection{Medien}
\begin{itemize}
    \item Tafel/Whiteboard
    \item Lehrwerk: \lehrwerk, \lehrwerkseite
    \item Arbeitsblatt: AB-06d.pdf (Wiederholung)
    \item Stundenleistung: SL-06d.pdf (Test, 20P)
    \item Lernpfad: LP-06d.html (interaktive Wiederholung)
    \item Stadtplan-Grafik (für Wegbeschreibung)
\end{itemize}

\subsection{Sozialformen}
\begin{itemize}
    \item Plenum: Wiederholung, Fragen klären
    \item Partnerarbeit: Dialog-Übung Wegbeschreibung
    \item Einzelarbeit: Stundenleistung
\end{itemize}

\subsection{Besonderheit: Stundenleistung}
\begin{itemize}
    \item Dauer: 25-30 Minuten
    \item Umfang: 20 Punkte (4 Aufgaben)
    \item Inhalt: Orte, ir, estar, Wegbeschreibung
    \item Bewertung: Note nach Prozentskala
    \item Differenzierung: Basis 60\% = ausreichend
\end{itemize}

% ============================================================================
% 6. VERLAUFSPLANUNG
% ============================================================================
\section{Verlaufsplanung}

\textbf{1. Stunde (45 Min.) -- Schwerpunkt: Wiederholung}

\begin{longtable}{|p{1cm}|p{1.5cm}|p{4cm}|p{3cm}|p{2cm}|p{2cm}|}
\hline
\textbf{Zeit} & \textbf{Phase} & \textbf{Lehrerhandeln} & \textbf{Schülerhandeln} & \textbf{SF} & \textbf{Medien} \\
\hline
\endhead

0--5 & Einstieg & ``¿Adónde vas hoy?'' Blitzlicht & Orte nennen & Plenum & -- \\
\hline

5--15 & Wiederholung 1 & Orte-Memory: Bild-Wort-Zuordnung & Zuordnen, Nennen & Plenum & Tafel \\
\hline

15--25 & Wiederholung 2 & ir + estar: Konjugation wiederholen & Tabelle ergänzen & EA & AB \\
\hline

25--35 & Übung & PA: Dialog Wegbeschreibung üben & Fragen und beschreiben & PA & AB \\
\hline

35--45 & Sicherung & LP-06d: Quiz zur Selbstüberprüfung & Bearbeiten & EA & LP \\
\hline
\end{longtable}

\vspace{1em}

\textbf{2. Stunde (45 Min.) -- Schwerpunkt: Stundenleistung}

\begin{longtable}{|p{1cm}|p{1.5cm}|p{4cm}|p{3cm}|p{2cm}|p{2cm}|}
\hline
\textbf{Zeit} & \textbf{Phase} & \textbf{Lehrerhandeln} & \textbf{Schülerhandeln} & \textbf{SF} & \textbf{Medien} \\
\hline
\endhead

0--5 & Aktivierung & Letzte Fragen klären & Fragen stellen & Plenum & -- \\
\hline

5--10 & Vorbereitung & Stundenleistung erklären & Zuhören, Fragen & Plenum & SL \\
\hline

10--35 & \textbf{TEST} & Aufsicht, individuelle Hilfe bei Verständnisfragen & Stundenleistung bearbeiten & EA & SL \\
\hline

35--40 & Abgabe & Einsammeln, kurze Entspannung & Abgeben & -- & -- \\
\hline

40--45 & Abschluss & Ausblick Sequenz 7, Feedback & Rückmeldung & Plenum & -- \\
\hline
\end{longtable}

\textbf{Legende:} EA = Einzelarbeit, PA = Partnerarbeit, SF = Sozialform, SL = Stundenleistung, LP = Lernpfad

% ============================================================================
% 7. ERWARTETE SCHWIERIGKEITEN
% ============================================================================
\section{Erwartete Schwierigkeiten}

\begin{tabular}{|p{5cm}|p{8cm}|}
\hline
\textbf{Schwierigkeit} & \textbf{Reaktion} \\
\hline
Verwechslung ir/estar & Bedeutungsunterschied: ir = gehen ZU, estar = sein AN einem Ort \\
\hline
a + el nicht zu ``al'' kontrahiert & Regel explizit üben: a + el = al (aber: a + la bleibt) \\
\hline
Wegbeschreibungen unsicher & Gesten verwenden, Stadtplan zeigen \\
\hline
Prüfungsangst bei Stundenleistung & Ruhige Atmosphäre, klare Instruktionen, Zeitzugabe für \B{} \\
\hline
\end{tabular}

% ============================================================================
% 8. HAUSAUFGABE
% ============================================================================
\section{Hausaufgabe}

\begin{itemize}
    \item \textbf{Vor der Stunde:} LP-06d.html zur Wiederholung (optional)
    \item \textbf{Nach der Stunde:} Vokabeln Sequenz 6 wiederholen
\end{itemize}

% ============================================================================
% 9. LÖSUNGEN
% ============================================================================
\section{Lösungen AB-06d}

\textbf{Aufgabe 1 (Orte zuordnen):} (4P)
\begin{itemize}
    \item el parque = der Park
    \item la tienda = das Geschäft
    \item el cine = das Kino
    \item la farmacia = die Apotheke
\end{itemize}

\textbf{Aufgabe 2 (ir konjugieren):} (4P)
\begin{itemize}
    \item yo voy, tú vas, él/ella va, nosotros vamos
\end{itemize}

\textbf{Aufgabe 3 (estar + Ortsangaben):} (6P)
\begin{itemize}
    \item El banco está cerca del parque.
    \item La farmacia está enfrente del hospital.
    \item El museo está al lado de la iglesia.
\end{itemize}

\textbf{Aufgabe 4 (Wegbeschreibung):} (6P)
\begin{itemize}
    \item Sigue todo recto.
    \item Gira a la derecha.
    \item El cine está a la izquierda.
\end{itemize}

% ============================================================================
% 10. STUNDENLEISTUNG - ÜBERSICHT
% ============================================================================
\section{Stundenleistung SL-06d}

\textbf{Aufbau:} 4 Aufgaben, 20 Punkte, 25 Minuten

\begin{tabular}{|c|l|c|c|c|}
\hline
\textbf{Aufg.} & \textbf{Inhalt} & \textbf{AFB} & \textbf{Punkte} & \textbf{Diff.} \\
\hline
1 & Zuordnung: Orte Spanisch-Deutsch & I & 4 & alle \\
\hline
2 & Konjugation: ir und estar & I/II & 4 & alle \\
\hline
3 & Lückentext: Wegbeschreibung & II & 6 & alle \\
\hline
4 & Stadtplan: Wegbeschreibung schreiben & II/III & 6 & alle \\
\hline
\multicolumn{3}{|r|}{\textbf{Gesamt:}} & \textbf{20} & \\
\hline
\end{tabular}

\textbf{Bewertung (Klasse 7):}
\begin{tabular}{|c|c|c|c|c|c|c|}
\hline
Note & 1 & 2 & 3 & 4 & 5 & 6 \\
\hline
ab \% & 96 & 80 & 60 & 40 & 20 & <20 \\
\hline
ab Punkte & 19 & 16 & 12 & 8 & 4 & <4 \\
\hline
\end{tabular}

% ============================================================================
% 11. DIDAKTIK-SCORE
% ============================================================================
\newpage
\section{Didaktik-Score}

\begin{scorebox}
\textbf{Bewertung der didaktischen Qualität nach 10 Dimensionen}\\
\textbf{Ziel-Score: $\geq$ 9.0 (Premium-Qualität)}

\vspace{1em}
\begin{tabular}{|c|l|c|c|p{4.5cm}|}
\hline
\textbf{\#} & \textbf{Dimension} & \textbf{Gew.} & \textbf{Score} & \textbf{Notiz} \\
\hline
1 & Differenzierung & 15\% & 9/10 & SL mit Basis-Sicherung, Zusatzaufgaben \\
\hline
2 & Taxonomie (AFB) & 10\% & 9/10 & AFB I: 40\%, AFB II: 45\%, AFB III: 15\% \\
\hline
3 & Lernziel-Alignment & 10\% & 10/10 & Rahmenplan: Wegbeschreibung \\
\hline
4 & Aktivierung & 10\% & 9/10 & Memory, Dialog-PA, Lernpfad \\
\hline
5 & Scaffolding & 10\% & 9/10 & Wiederholung $\rightarrow$ Übung $\rightarrow$ Test \\
\hline
6 & Feedback-Qualität & 10\% & 9/10 & LP mit Sofortfeedback, Bewertungsraster \\
\hline
7 & Sprachliche Klarheit & 10\% & 10/10 & Altersgerecht Kl. 7 \\
\hline
8 & Motivation \& Relevanz & 10\% & 9/10 & Alltagsrelevanz: Orientierung in Stadt \\
\hline
9 & Inklusion (UDL) & 10\% & 9/10 & Stadtplan visuell, Zeitzugabe \\
\hline
10 & Praktikabilität & 5\% & 9/10 & 45+45 Min realistisch \\
\hline
\multicolumn{3}{|r|}{\textbf{GESAMT:}} & \textbf{9.2/10} & \textcolor{successgreen}{\textbf{FREIGABE}} \\
\hline
\end{tabular}
\end{scorebox}

\end{document}
